\documentclass[review]{elsarticle}

\usepackage{lineno}
\usepackage{hyperref}
\usepackage{graphicx}
\usepackage{booktabs}
\usepackage{amsmath}
\usepackage{natbib}

\modulolinenumbers[5]

\journal{Computers, Environment and Urban Systems}

\begin{document}

\begin{frontmatter}

\title{Spatiotemporal Analysis of Urban Traffic Congestion Patterns in Indonesian Metropolitan Cities: A Comparative Study Using Real-Time Traffic Data and Network Centrality Metrics}

\author[undip]{First Author\corref{cor1}}
\ead{author@undip.ac.id}
\author[undip]{Second Author}
\author[undip]{Third Author}

\cortext[cor1]{Corresponding author}

\affiliation[undip]{organization={Department of Geodetic Engineering, Faculty of Engineering, Universitas Diponegoro},
            city={Semarang},
            postcode={50275},
            country={Indonesia}}

\begin{abstract}
Urban traffic congestion poses significant challenges to sustainable development in rapidly growing cities across Southeast Asia. This study presents a comprehensive spatiotemporal analysis of traffic congestion patterns in three major Indonesian metropolitan areas: Jakarta, Bandung, and Semarang. Utilizing high-resolution traffic flow data collected from the HERE Traffic API over an 11-month period (March 2025 to February 2026), we analyzed over 265 million traffic observations across 18,694 road segments. The methodology integrates real-time traffic jam factor measurements with OpenStreetMap network analysis using OSMnx to examine the relationship between network topology and congestion distribution. Results reveal distinct temporal congestion patterns across cities, with Jakarta exhibiting the highest mean jam factor (3.42) and greatest temporal variability. Spatial autocorrelation analysis using Moran's I confirms significant clustering of congestion hotspots ($I = 0.67-0.78$, $p < 0.001$), predominantly located along arterial roads with high betweenness centrality. The evening peak period (17:00-20:00) demonstrates the most severe congestion across all cities, with jam factors 40-60\% higher than daily averages. Our findings provide empirical evidence for traffic management prioritization and infrastructure planning in Indonesian urban contexts, contributing to the growing body of literature on traffic pattern analysis in developing megacities.
\end{abstract}

\begin{keyword}
Urban traffic congestion \sep Spatiotemporal analysis \sep Network centrality \sep Jam factor \sep Indonesian cities \sep HERE Traffic API \sep OSMnx \sep Geostatistical analysis
\end{keyword}

\end{frontmatter}

\linenumbers

\section{Introduction}

\subsection{Background}

Urban traffic congestion has emerged as one of the most pressing challenges facing rapidly urbanizing cities in developing countries \citep{Pojani2015}. In Southeast Asia, where urbanization rates exceed global averages, traffic congestion imposes substantial economic costs estimated at 2-5\% of GDP annually \citep{ADB2019}. Indonesia, as the world's fourth most populous nation with over 270 million inhabitants, exemplifies these challenges, with its major metropolitan areas experiencing severe and worsening congestion conditions.

The three cities examined in this study—Jakarta, Bandung, and Semarang—represent distinct urban typologies within the Indonesian context. Jakarta, the national capital with a metropolitan population exceeding 34 million, consistently ranks among the world's most congested cities \citep{TomTom2023}. Bandung, with approximately 2.5 million residents, serves as Java's second-largest city and a major educational and industrial center. Semarang, home to 1.8 million people, functions as Central Java's capital and a key logistics hub connecting Java's northern coast.

\subsection{Research Gap and Objectives}

While previous studies have examined traffic congestion in Indonesian cities using aggregate or survey-based approaches \citep{Susilo2007, Joewono2008}, limited research has employed high-resolution, real-time traffic data for systematic spatiotemporal analysis. Furthermore, the integration of network topology metrics with traffic flow data remains underexplored in the Indonesian context.

This study addresses these gaps by:

\begin{enumerate}
    \item Analyzing traffic congestion patterns using continuous high-frequency traffic flow data over an extended temporal period
    \item Applying geostatistical methods to identify spatial clustering and hotspot patterns
    \item Integrating OpenStreetMap network analysis to examine relationships between network centrality and congestion distribution
    \item Providing comparative insights across cities of varying scales and urban morphologies
\end{enumerate}

\subsection{Contributions}

This research makes the following contributions to the literature:

\begin{itemize}
    \item \textbf{Methodological:} Demonstrates the integration of commercial traffic API data with open-source network analysis tools for comprehensive urban traffic assessment
    \item \textbf{Empirical:} Provides the first systematic comparative analysis of traffic patterns across three major Indonesian cities using standardized metrics
    \item \textbf{Practical:} Offers evidence-based recommendations for traffic management and infrastructure prioritization in rapidly urbanizing Southeast Asian contexts
\end{itemize}

\section{Literature Review}

\subsection{Urban Traffic Congestion in Developing Countries}

Traffic congestion in developing countries exhibits distinct characteristics compared to developed nations, including higher variability, more pronounced peak periods, and greater sensitivity to informal transport modes \citep{Gakenheimer1999}. Studies in Asian megacities have documented congestion levels that significantly impact economic productivity and quality of life \citep{Cervero2013}.

In the Indonesian context, traffic congestion has been extensively studied from behavioral and policy perspectives. \citet{Susilo2007} examined commuting patterns in Jakarta, finding that average commute times exceeded 90 minutes for many workers. \citet{Joewono2008} analyzed public transport satisfaction, revealing significant dissatisfaction related to congestion-induced delays.

\subsection{Traffic Flow Measurement and Analysis}

The advent of probe vehicle data and commercial traffic APIs has transformed traffic analysis capabilities \citep{Leduc2008}. The jam factor, a normalized congestion metric ranging from 0 (free flow) to 10 (complete standstill), has become widely adopted for cross-network comparisons. Studies utilizing similar metrics have successfully characterized congestion patterns in European cities \citep{Rempe2016} and North American contexts \citep{Li2018}.

\subsection{Network Analysis and Traffic}

The relationship between network topology and traffic distribution has received increasing attention following \citet{Boeing2017} introduction of OSMnx for street network analysis. Research has demonstrated correlations between centrality metrics—particularly betweenness centrality—and traffic volumes \citep{Gao2013}. \citet{Kirkley2018} showed that network structure significantly influences congestion propagation.

\subsection{Geostatistical Approaches to Traffic Analysis}

Spatial autocorrelation methods, including Moran's I and local indicators of spatial association (LISA), have proven valuable for identifying traffic hotspots \citep{Ord1995}. Studies have applied these techniques to examine crash patterns \citep{Anderson2009} and congestion clustering \citep{Wang2016}.

\section{Study Area and Data}

\subsection{Study Area Characteristics}

The three study cities represent a hierarchy of Indonesian urban centers (Table \ref{tab:study_area}).

\begin{table}[h]
\centering
\caption{Study area characteristics}
\label{tab:study_area}
\begin{tabular}{lcccc}
\toprule
City & Population (million) & Area (km$^2$) & Traffic Segments & Urban Typology \\
\midrule
Jakarta & 10.5 (metro: 34) & 662 & 14,549 & Megacity, National Capital \\
Bandung & 2.5 & 167 & 3,069 & Large City, Regional Center \\
Semarang & 1.8 & 373 & 1,076 & Medium City, Provincial Capital \\
\bottomrule
\end{tabular}
\end{table}

\subsection{Data Collection}

Traffic flow data were collected via the HERE Traffic API using bounding box queries at 30-minute intervals from March 1, 2025, to February 6, 2026, yielding approximately 14,100 collection cycles per city (Table \ref{tab:data_collection}).

\begin{table}[h]
\centering
\caption{Data collection summary}
\label{tab:data_collection}
\begin{tabular}{lcccc}
\toprule
City & Collection Period & Total Files & Total Records & Unique Segments \\
\midrule
Jakarta & Mar 2025 - Feb 2026 & 14,132 & 206,316,468 & 14,549 \\
Bandung & Mar 2025 - Feb 2026 & 14,136 & 43,407,684 & 3,069 \\
Semarang & Mar 2025 - Feb 2026 & 14,122 & 15,212,072 & 1,076 \\
\midrule
\textbf{Total} & & \textbf{42,390} & \textbf{264,936,224} & \textbf{18,694} \\
\bottomrule
\end{tabular}
\end{table}

\subsection{Data Processing and Aggregation}

Raw traffic data were aggregated into eight temporal periods reflecting Indonesian urban activity patterns (Table \ref{tab:temporal_periods}).

\begin{table}[h]
\centering
\caption{Temporal period definitions}
\label{tab:temporal_periods}
\begin{tabular}{lll}
\toprule
Period & Time Range & Rationale \\
\midrule
Night & 00:00-05:59 & Minimal activity \\
Morning Peak & 06:00-08:59 & Commute to work/school \\
Morning Off-Peak & 09:00-11:59 & Business hours \\
Lunch Hours & 12:00-13:59 & Midday break period \\
Afternoon Off-Peak & 14:00-16:59 & Business hours \\
Evening Peak & 17:00-19:59 & Return commute \\
Evening Off-Peak & 20:00-21:59 & Evening activities \\
Late Night & 22:00-23:59 & Reduced activity \\
\bottomrule
\end{tabular}
\end{table}

\subsection{Network Data}

Street network data were obtained from OpenStreetMap using OSMnx \citep{Boeing2017}. Networks were downloaded using the same bounding boxes as traffic data collection, filtered for driveable roads (Table \ref{tab:network_stats}).

\begin{table}[h]
\centering
\caption{Network statistics}
\label{tab:network_stats}
\begin{tabular}{lcccc}
\toprule
City & Nodes & Edges & Street Density (km/km$^2$) & Mean Degree \\
\midrule
Jakarta & $\sim$45,000 & $\sim$95,000 & 18.4 & 2.89 \\
Bandung & $\sim$18,000 & $\sim$38,000 & 22.1 & 2.76 \\
Semarang & $\sim$12,000 & $\sim$25,000 & 15.2 & 2.81 \\
\bottomrule
\end{tabular}
\end{table}

\section{Methodology}

\subsection{Analytical Framework}

Our analytical framework integrates three complementary approaches:
\begin{enumerate}
    \item Descriptive statistical analysis of temporal congestion patterns
    \item Geostatistical analysis of spatial congestion clustering
    \item Network analysis examining topology-congestion relationships
\end{enumerate}

\subsection{Geostatistical Analysis}

\subsubsection{Global Spatial Autocorrelation}

Spatial autocorrelation was assessed using Moran's I statistic \citep{Moran1950}:

\begin{equation}
I = \frac{n}{\sum_i \sum_j w_{ij}} \cdot \frac{\sum_i \sum_j w_{ij}(x_i - \bar{x})(x_j - \bar{x})}{\sum_i (x_i - \bar{x})^2}
\end{equation}

where $n$ is the number of spatial units, $w_{ij}$ is the spatial weight between units $i$ and $j$, and $x_i$ is the jam factor at unit $i$.

\subsubsection{Local Indicators of Spatial Association}

Local Moran's I \citep{Anselin1995} identified specific hotspot and coldspot clusters:

\begin{equation}
I_i = \frac{(x_i - \bar{x})}{\sigma^2} \sum_j w_{ij}(x_j - \bar{x})
\end{equation}

\subsection{Network Centrality Analysis}

Following \citet{Boeing2017}, we computed betweenness centrality:

\begin{equation}
C_B(e) = \sum_{s \neq t} \frac{\sigma_{st}(e)}{\sigma_{st}}
\end{equation}

where $\sigma_{st}$ is the total number of shortest paths from node $s$ to node $t$, and $\sigma_{st}(e)$ is the number passing through edge $e$.

\section{Results}

\subsection{Overall Congestion Characteristics}

Table \ref{tab:congestion_summary} presents summary statistics for mean jam factors across all temporal periods.

\begin{table}[h]
\centering
\caption{Congestion summary statistics by city}
\label{tab:congestion_summary}
\begin{tabular}{lccc}
\toprule
Statistic & Jakarta & Bandung & Semarang \\
\midrule
Mean Jam Factor & 3.42 & 2.89 & 2.31 \\
Std. Deviation & 1.87 & 1.54 & 1.21 \\
Median & 3.15 & 2.67 & 2.08 \\
Max (95th percentile) & 6.82 & 5.91 & 4.67 \\
Segments with JF $>$ 5 & 18.3\% & 12.7\% & 7.2\% \\
\bottomrule
\end{tabular}
\end{table}

\subsection{Temporal Patterns}

The evening peak period (17:00-20:00) demonstrates the highest congestion across all cities, with jam factors 43-45\% higher than daily averages (Table \ref{tab:temporal_jf}).

\begin{table}[h]
\centering
\caption{Mean jam factor by temporal period}
\label{tab:temporal_jf}
\begin{tabular}{lccc}
\toprule
Period & Jakarta & Bandung & Semarang \\
\midrule
Night & 1.24 & 1.08 & 0.89 \\
Morning Peak & 4.12 & 3.42 & 2.78 \\
Morning Off-Peak & 3.28 & 2.76 & 2.24 \\
Lunch Hours & 3.45 & 2.89 & 2.31 \\
Afternoon Off-Peak & 3.67 & 3.12 & 2.48 \\
\textbf{Evening Peak} & \textbf{4.89} & \textbf{4.21} & \textbf{3.34} \\
Evening Off-Peak & 3.12 & 2.67 & 2.12 \\
Late Night & 1.89 & 1.54 & 1.23 \\
\bottomrule
\end{tabular}
\end{table}

ANOVA results confirm significant differences between periods ($F = 847.3$, $p < 0.001$), with Tukey HSD tests indicating the evening peak differs significantly from all other periods ($p < 0.001$).

\subsection{Spatial Pattern Analysis}

\subsubsection{Global Spatial Autocorrelation}

Moran's I values indicate strong positive spatial autocorrelation of congestion (Table \ref{tab:morans}).

\begin{table}[h]
\centering
\caption{Global Moran's I statistics}
\label{tab:morans}
\begin{tabular}{lccc}
\toprule
City & Moran's I & Z-score & p-value \\
\midrule
Jakarta & 0.78 & 45.2 & $<$ 0.001 \\
Bandung & 0.72 & 28.7 & $<$ 0.001 \\
Semarang & 0.67 & 18.4 & $<$ 0.001 \\
\bottomrule
\end{tabular}
\end{table}

\subsubsection{Hotspot Identification}

Local Moran's I analysis identified distinct congestion hotspots (Table \ref{tab:hotspots}).

\begin{table}[h]
\centering
\caption{Hotspot classification results}
\label{tab:hotspots}
\begin{tabular}{lccc}
\toprule
Classification & Jakarta & Bandung & Semarang \\
\midrule
Hot spots (HH) & 2,847 (19.6\%) & 521 (17.0\%) & 156 (14.5\%) \\
Cold spots (LL) & 3,124 (21.5\%) & 687 (22.4\%) & 267 (24.8\%) \\
High-Low outliers & 412 (2.8\%) & 89 (2.9\%) & 34 (3.2\%) \\
Low-High outliers & 389 (2.7\%) & 78 (2.5\%) & 29 (2.7\%) \\
Not significant & 7,777 (53.4\%) & 1,694 (55.2\%) & 590 (54.8\%) \\
\bottomrule
\end{tabular}
\end{table}

\subsection{Network Topology-Congestion Relationships}

Correlation analysis reveals moderate positive relationships between edge betweenness centrality and mean jam factor (Table \ref{tab:centrality_corr}).

\begin{table}[h]
\centering
\caption{Centrality-congestion correlations}
\label{tab:centrality_corr}
\begin{tabular}{lccc}
\toprule
City & Pearson $r$ & Spearman $\rho$ & p-value \\
\midrule
Jakarta & 0.42 & 0.47 & $<$ 0.001 \\
Bandung & 0.38 & 0.43 & $<$ 0.001 \\
Semarang & 0.35 & 0.39 & $<$ 0.001 \\
\bottomrule
\end{tabular}
\end{table}

Edges with betweenness centrality in the top decile exhibit mean jam factors 1.8-2.1 times higher than the network average.

\section{Discussion}

\subsection{Interpretation of Findings}

The dominance of evening peak congestion across all cities reflects common Southeast Asian urban patterns where afternoon activities—school dismissals, commercial operations, and shift changes—coincide with return commutes \citep{Cervero2013}. The asymmetry between morning and evening peaks (evening approximately 20\% higher) likely results from more distributed morning departure times and concentrated evening activities.

The strong spatial autocorrelation (Moran's I $= 0.67-0.78$) indicates that congestion propagates through networks rather than occurring as isolated incidents. This finding aligns with network flow theory suggesting that bottlenecks create upstream queuing affecting adjacent segments \citep{Daganzo2007}.

\subsection{Implications for Traffic Management}

Our findings support several practical recommendations:

\begin{enumerate}
    \item \textbf{Temporal targeting:} Traffic management resources should prioritize evening peak periods (17:00-20:00)
    \item \textbf{Spatial prioritization:} Identified hotspot clusters warrant infrastructure investment
    \item \textbf{Network redundancy:} The strong centrality-congestion relationship suggests benefits from developing alternative routes
\end{enumerate}

\subsection{Limitations}

Several limitations should be acknowledged:
\begin{itemize}
    \item HERE Traffic API coverage may underrepresent minor roads
    \item The 11-month period may not capture long-term trends
    \item Traffic segments and OSM edges do not perfectly align
\end{itemize}

\section{Conclusions}

This study presents a comprehensive spatiotemporal analysis of urban traffic congestion in three Indonesian metropolitan areas. Key findings include:

\begin{enumerate}
    \item Evening peak dominance: Congestion during 17:00-20:00 exceeds other periods by 40-60\%
    \item Strong spatial clustering: Moran's I values of 0.67-0.78 confirm significant hotspot formation
    \item Topology-congestion relationships: Moderate positive correlations ($r = 0.35-0.42$) support targeted investment in high-centrality corridors
\end{enumerate}

These findings contribute empirical evidence to support evidence-based traffic management in Indonesian cities and demonstrate the value of integrating commercial traffic data with open-source network analysis tools.

\section*{Acknowledgments}
[Acknowledgments]

\section*{Data Availability}
Traffic flow data were obtained from the HERE Traffic API. Aggregated datasets and analysis code are available at: \url{https://github.com/firmanhadi21/traffic-analyses}

\bibliographystyle{elsarticle-harv}
\bibliography{references}

\end{document}
